\documentclass[12pt, a4paper]{article}

% ── Page geometry ──────────────────────────────────────────────────────────────
\usepackage[a4paper, top=2.5cm, bottom=2.5cm, left=2.8cm, right=2.8cm,
            headheight=15pt]{geometry}

% ── Font & encoding ────────────────────────────────────────────────────────────
\usepackage[T1]{fontenc}
\usepackage[utf8]{inputenc}
\usepackage{palatino}
\usepackage{mathpazo}
\usepackage{microtype}

% ── Maths ─────────────────────────────────────────────────────────────────────
\usepackage{amsmath, amssymb, amsthm, mathtools}
\usepackage{cancel}

% ── Graphics & colour ─────────────────────────────────────────────────────────
\usepackage{xcolor}
\usepackage{tikz}
\usetikzlibrary{arrows.meta, positioning, calc, shapes.geometric,
                decorations.pathreplacing, backgrounds, fit, patterns}
\usepackage{pgfplots}
\pgfplotsset{compat=1.15}

% ── Tables ────────────────────────────────────────────────────────────────────
\usepackage{booktabs, array, colortbl, multirow, tabularx}
\renewcommand{\arraystretch}{1.4}

% ── Lists ─────────────────────────────────────────────────────────────────────
\usepackage{enumitem}
\setlist[itemize]{leftmargin=1.8em, itemsep=3pt, topsep=4pt}
\setlist[enumerate]{leftmargin=1.8em, itemsep=3pt, topsep=4pt}

% ── Boxes ─────────────────────────────────────────────────────────────────────
\usepackage[most]{tcolorbox}
\tcbuselibrary{skins, breakable, theorems}

% ── Headers & footers ─────────────────────────────────────────────────────────
\usepackage{fancyhdr}
\usepackage{lastpage}

% ── Hyperlinks ────────────────────────────────────────────────────────────────
\usepackage[colorlinks=true, linkcolor=royalblue, urlcolor=royalblue,
            citecolor=royalblue, bookmarks=true]{hyperref}
\usepackage{tocloft}

% ── Miscellaneous ─────────────────────────────────────────────────────────────
\usepackage{parskip}
\setlength{\parskip}{6pt}
\setlength{\parindent}{0pt}
\usepackage{float}
\usepackage{caption}

% ═══════════════════════════════════════════════════════════════════════════════
%  COLOUR PALETTE  (deep navy / gold / teal)
% ═══════════════════════════════════════════════════════════════════════════════
\definecolor{royalblue}{RGB}{25, 55, 130}
\definecolor{navylight}{RGB}{235, 239, 252}
\definecolor{navyborder}{RGB}{100, 130, 200}
\definecolor{gold}{RGB}{185, 135, 15}
\definecolor{lightgold}{RGB}{254, 248, 225}
\definecolor{teal}{RGB}{15, 115, 105}
\definecolor{lightteal}{RGB}{225, 248, 244}
\definecolor{tealborder}{RGB}{15, 115, 105}
\definecolor{crimson}{RGB}{155, 25, 35}
\definecolor{lightcrimson}{RGB}{255, 235, 235}
\definecolor{purple}{RGB}{88, 35, 140}
\definecolor{lightpurple}{RGB}{245, 238, 255}
\definecolor{charcoal}{RGB}{42, 42, 42}
\definecolor{slategray}{RGB}{80, 90, 105}
\definecolor{tabhead}{RGB}{25, 55, 130}
\definecolor{tabeven}{RGB}{235, 240, 252}
\definecolor{solgreen}{RGB}{12, 100, 55}
\definecolor{lightsol}{RGB}{228, 252, 238}

% ═══════════════════════════════════════════════════════════════════════════════
%  TCOLORBOX ENVIRONMENTS
% ═══════════════════════════════════════════════════════════════════════════════

\newtcolorbox{defbox}[1]{
  breakable, enhanced,
  colback=navylight, colframe=royalblue, coltitle=white,
  fonttitle=\bfseries\sffamily, title={Definition: #1},
  arc=4pt, boxrule=1pt, left=8pt, right=8pt, top=6pt, bottom=6pt,
  attach boxed title to top left={yshift=-2mm, xshift=6mm},
  boxed title style={colback=royalblue, arc=3pt, boxrule=0pt}
}

\newtcolorbox{thmbox}[1]{
  breakable, enhanced,
  colback=lightteal, colframe=teal, coltitle=white,
  fonttitle=\bfseries\sffamily, title={Theorem: #1},
  arc=4pt, boxrule=1pt, left=8pt, right=8pt, top=6pt, bottom=6pt,
  attach boxed title to top left={yshift=-2mm, xshift=6mm},
  boxed title style={colback=teal, arc=3pt, boxrule=0pt}
}

\newtcolorbox{propbox}[1]{
  breakable, enhanced,
  colback=lightteal, colframe=teal, coltitle=white,
  fonttitle=\bfseries\sffamily, title={Proposition: #1},
  arc=4pt, boxrule=1pt, left=8pt, right=8pt, top=6pt, bottom=6pt,
  attach boxed title to top left={yshift=-2mm, xshift=6mm},
  boxed title style={colback=teal, arc=3pt, boxrule=0pt}
}

\newtcolorbox{keybox}[1]{
  breakable, enhanced,
  colback=lightgold, colframe=gold, coltitle=charcoal,
  fonttitle=\bfseries\sffamily, title={#1},
  arc=4pt, boxrule=1pt, left=8pt, right=8pt, top=6pt, bottom=6pt,
  attach boxed title to top left={yshift=-2mm, xshift=6mm},
  boxed title style={colback=gold, arc=3pt, boxrule=0pt, coltitle=white}
}

\newtcolorbox{exbox}[1]{
  breakable, enhanced,
  colback=lightpurple, colframe=purple, coltitle=white,
  fonttitle=\bfseries\sffamily, title={Worked Example: #1},
  arc=4pt, boxrule=1pt, left=8pt, right=8pt, top=6pt, bottom=6pt,
  attach boxed title to top left={yshift=-2mm, xshift=6mm},
  boxed title style={colback=purple, arc=3pt, boxrule=0pt}
}

\newtcolorbox{warnbox}[1]{
  breakable, enhanced,
  colback=lightcrimson, colframe=crimson, coltitle=white,
  fonttitle=\bfseries\sffamily, title={#1},
  arc=4pt, boxrule=1pt, left=8pt, right=8pt, top=6pt, bottom=6pt,
  attach boxed title to top left={yshift=-2mm, xshift=6mm},
  boxed title style={colback=crimson, arc=3pt, boxrule=0pt}
}

\newtcolorbox{solbox}{
  breakable, enhanced,
  colback=lightsol, colframe=solgreen, coltitle=white,
  fonttitle=\bfseries\sffamily, title={Solution},
  arc=4pt, boxrule=1pt, left=8pt, right=8pt, top=6pt, bottom=6pt,
  attach boxed title to top left={yshift=-2mm, xshift=6mm},
  boxed title style={colback=solgreen, arc=3pt, boxrule=0pt}
}

% ── Theorem-style environments ─────────────────────────────────────────────────
\theoremstyle{definition}
\newtheorem{remark}{Remark}[section]
\newtheorem{corollary}{Corollary}[section]

% ═══════════════════════════════════════════════════════════════════════════════
%  MATHS MACROS
% ═══════════════════════════════════════════════════════════════════════════════
\newcommand{\R}{\mathbb{R}}
\newcommand{\N}{\mathbb{N}}
\newcommand{\Z}{\mathbb{Z}}
\newcommand{\Q}{\mathbb{Q}}
\newcommand{\C}{\mathbb{C}}
\newcommand{\Rn}{\mathbb{R}^n}
\newcommand{\Rm}{\mathbb{R}^m}
\newcommand{\powerset}{\mathcal{P}}
\newcommand{\st}{\;\text{s.t.}\;}
\newcommand{\xto}[1]{\xrightarrow{#1}}
\newcommand{\inv}{^{-1}}
\DeclareMathOperator{\dom}{dom}
\DeclareMathOperator{\ran}{ran}
\DeclareMathOperator{\img}{img}

% ── Table helpers ──────────────────────────────────────────────────────────────
\newcommand{\thead}[1]{\cellcolor{tabhead}\color{white}\textbf{#1}}

% ═══════════════════════════════════════════════════════════════════════════════
%  PAGE STYLE
% ═══════════════════════════════════════════════════════════════════════════════
\pagestyle{fancy}
\fancyhf{}
\fancyhead[L]{\textcolor{royalblue}{\textbf{ECON 320}} \textcolor{slategray}{--- Sets \& Functions}}
\fancyhead[R]{\textcolor{slategray}{\small School of Economics}}
\fancyfoot[C]{\textcolor{slategray}{\small Page \thepage\ of \pageref{LastPage}}}
\renewcommand{\headrulewidth}{1.2pt}
\renewcommand{\headrule}{%
  \hbox to\headwidth{\color{royalblue}\leaders\hrule height\headrulewidth\hfill}}
\renewcommand{\footrulewidth}{0.4pt}

% ── TOC styling ───────────────────────────────────────────────────────────────
\renewcommand{\cftsecfont}{\bfseries\color{royalblue}}
\renewcommand{\cftsecpagefont}{\bfseries\color{royalblue}}
\renewcommand{\cftsubsecfont}{\color{charcoal}}
\setlength{\cftbeforesecskip}{4pt}

% ═══════════════════════════════════════════════════════════════════════════════
\begin{document}
% ═══════════════════════════════════════════════════════════════════════════════

% ── COVER ─────────────────────────────────────────────────────────────────────
\thispagestyle{empty}
\begin{tikzpicture}[remember picture, overlay]
  \fill[royalblue] (current page.north west)
    rectangle ([yshift=-5.8cm] current page.north east);
  \fill[gold] ([yshift=-5.8cm] current page.north west)
    rectangle ([yshift=-6.2cm] current page.north east);
\end{tikzpicture}

\vspace*{-1.5cm}
\begin{center}
  {\color{white}\large\bfseries\sffamily ECON 320 -- Solutions to Practice Problem Set 1}\\[10pt]
  {\color{white}\Huge\bfseries Sets and Functions}\\[8pt]
  {\color{gold}\large\itshape Mathematical Foundations for Economic Analysis}
\end{center}

\vspace{3.4cm}

\begin{center}
\begin{tabular}{l l}
  \textbf{\textcolor{royalblue}{Course:}}        & ECON 320: Introduction to Mathematical Economics \\[3pt]


  \textbf{\textcolor{royalblue}{Term:}}          & Fall 2026 \\[3pt]
    \textbf{\textcolor{royalblue}{Instructor:}}          & Muhebullah Karimzada \\[3pt]
\end{tabular}
\end{center}




\newpage





\newpage
% ═══════════════════════════════════════════════════════════════════════════════
\section{Full Solutions}
% ═══════════════════════════════════════════════════════════════════════════════

\subsection{Solutions: Section A}

\begin{solbox}
\textbf{A1. Set Builder Notation}

\textbf{(a)} Even natural numbers less than 20:
\[A = \{x \in \N : x \text{ is even and } x < 20\} = \{2, 4, 6, 8, 10, 12, 14, 16, 18\}.\quad\text{Finite (9 elements).}\]

\textbf{(b)} Real numbers with $x^2 < 16$:
\[B = \{x \in \R : x^2 < 16\} = \{x \in \R : |x| < 4\} = (-4, 4).\quad\text{Infinite.}\]

\textbf{(c)} Prices with $D(P) = 50 - P > 0$ and $P \in [0,100]$:

$D(P) > 0 \Rightarrow P < 50$. Combined with $P \in [0,100]$:
\[C = \{P \in [0,100] : 50 - P > 0\} = [0, 50).\quad\text{Infinite.}\]
\end{solbox}

\begin{solbox}
\textbf{A2. Set Operations}

With $U=\{1,\ldots,10\}$, $A=\{1,2,3,4,5\}$, $B=\{3,4,5,6,7\}$, $C=\{5,6,7,8,9\}$:

\textbf{(a)} $A \cup B = \{1,2,3,4,5,6,7\}$

\textbf{(b)} $A \cap B = \{3,4,5\}$

\textbf{(c)} $A \setminus B = \{1,2\}$\quad (elements in $A$ not in $B$)

\textbf{(d)} $B \setminus A = \{6,7\}$\quad (elements in $B$ not in $A$)

\textbf{(e)} $(A \cup B)^c = U \setminus \{1,2,3,4,5,6,7\} = \{8,9,10\}$

\textbf{(f)} $A \cap B \cap C = \{3,4,5\} \cap \{5,6,7,8,9\} = \{5\}$

\textbf{(g)} $B \cap C = \{5,6,7\}$, so $A \cup (B \cap C) = \{1,2,3,4,5\} \cup \{5,6,7\} = \{1,2,3,4,5,6,7\}$

\textbf{Check distributivity}: $A \cup (B \cap C) = (A \cup B) \cap (A \cup C)$.
$A \cup B = \{1,2,3,4,5,6,7\}$, $A \cup C = \{1,2,3,4,5,6,7,8,9\}$.
$(A\cup B)\cap(A\cup C) = \{1,2,3,4,5,6,7\}$. \checkmark
\end{solbox}

\begin{solbox}
\textbf{A3. Power Set}

$A = \{x, y\}$, so $|A|=2$ and $|\powerset(A)| = 2^2 = 4$:
\[\powerset(A) = \{\emptyset, \{x\}, \{y\}, \{x,y\}\}\]
Now $|\powerset(A)| = 4$, so $|\powerset(\powerset(A))| = 2^4 = \mathbf{16}$.
\end{solbox}
\newpage
\begin{solbox}
\textbf{A4. Cartesian Product}

$A\times B = \{(1,a),(1,b),(1,c),(2,a),(2,b),(2,c)\}$\quad ($|A\times B| = 2\times 3 = 6$ elements)

$B\times A = \{(a,1),(a,2),(b,1),(b,2),(c,1),(c,2)\}$\quad ($|B\times A| = 3\times 2 = 6$ elements)

$A\times B \neq B\times A$: the pairs are ordered differently; e.g.\ $(1,a)\in A\times B$ but $(1,a)\notin B\times A$ since $1\notin B$. In general $A\times B = B\times A$ only if $A = B$.
\end{solbox}

\begin{solbox}
\textbf{A5. De Morgan's Law Verification}

$U=\{1,\ldots,8\}$, $A=\{1,2,3,4\}$, $B=\{3,4,5,6\}$.

\textbf{Left side}: $A\cup B = \{1,2,3,4,5,6\}$, so $(A\cup B)^c = \{7,8\}$.

\textbf{Right side}: $A^c = \{5,6,7,8\}$, $B^c = \{1,2,7,8\}$, so $A^c\cap B^c = \{7,8\}$.

$(A\cup B)^c = A^c \cap B^c = \{7,8\}$. \checkmark
\end{solbox}

\begin{solbox}
\textbf{A6. Budget Set Intersection}

$A = \{(x_1,x_2)\in\R^2_+ : 2x_1+5x_2\leq 100\}$ is the full budget set (triangle with vertices $(0,0)$, $(50,0)$, $(0,20)$).

$B = \{(x_1,x_2)\in\R^2_+ : x_1 \geq 10\}$ is the closed half-plane to the right of $x_1=10$.

$A\cap B = \{(x_1,x_2)\in\R^2_+ : 2x_1+5x_2\leq 100,\ x_1\geq 10\}$

Geometrically: the budget triangle \emph{truncated} by the vertical line $x_1=10$. It is a quadrilateral with vertices $(10,0)$, $(50,0)$, $(0,20)$... wait --- at $x_1=10$: $2(10)+5x_2\leq 100 \Rightarrow x_2\leq 16$. So the binding vertices are $(10,0)$, $(50,0)$, and $(10,16)$.

\textbf{Economic interpretation}: affordable bundles that contain \emph{at least 10 units of good 1} (perhaps a minimum consumption requirement or a rationing constraint). This is the feasible set when the consumer is forced to consume at least 10 units of good 1.
\end{solbox}

\begin{solbox}
\textbf{A7. Proof of Distributive Law}

We prove $A\cap(B\cup C) = (A\cap B)\cup(A\cap C)$ by element-chasing.

\textbf{($\subseteq$)}: Let $x \in A\cap(B\cup C)$. Then $x\in A$ and $x\in B\cup C$.

\textit{Case 1}: $x\in B$. Then $x\in A$ and $x\in B$, so $x\in A\cap B\subseteq (A\cap B)\cup(A\cap C)$.

\textit{Case 2}: $x\in C$. Then $x\in A$ and $x\in C$, so $x\in A\cap C\subseteq (A\cap B)\cup(A\cap C)$.

In either case, $x\in(A\cap B)\cup(A\cap C)$.

\textbf{($\supseteq$)}: Let $x\in(A\cap B)\cup(A\cap C)$.

\textit{Case 1}: $x\in A\cap B$. Then $x\in A$ and $x\in B\subseteq B\cup C$, so $x\in A\cap(B\cup C)$.

\textit{Case 2}: $x\in A\cap C$. Then $x\in A$ and $x\in C\subseteq B\cup C$, so $x\in A\cap(B\cup C)$.

Since $A\cap(B\cup C)\subseteq(A\cap B)\cup(A\cap C)$ and $(A\cap B)\cup(A\cap C)\subseteq A\cap(B\cup C)$, equality holds. \qed
\end{solbox}

\begin{solbox}
\textbf{A8. Supremum and Infimum}

\textbf{(a)} $A=\left\{\frac{n}{n+1}:n\in\N\right\} = \left\{\frac{1}{2},\frac{2}{3},\frac{3}{4},\ldots\right\}$

The sequence $n/(n+1)$ is strictly increasing and approaches 1. So $\sup A = 1$ (not attained: $n/(n+1)<1$ for all $n$, so $\max A$ does not exist). $\inf A = \min A = \frac{1}{2}$ (attained at $n=1$).

\textbf{(b)} $B=(-3,5]$: $\inf B = -3$ (not attained, $\min B$ DNE); $\sup B = \max B = 5$.

\textbf{(c)} $C=\{x\in\R:x^2<9\} = (-3,3)$: $\inf C = -3$ (DNE); $\sup C = 3$ (DNE). Neither $\min$ nor $\max$ exists.

\textbf{(d)} $D=\{x\in\R:x^3\leq 8\} = (-\infty,2]$: $\sup D = \max D = 2$. $D$ is unbounded below: $\inf D = -\infty$ (no infimum in $\R$).
\end{solbox}

\subsection{Solutions: Section B}

\begin{solbox}
\textbf{B1. Domain and Range}

\textbf{(a)} $f(x)=\sqrt{9-x^2}$: need $9-x^2\geq 0$, i.e.\ $x^2\leq 9$, i.e.\ $x\in[-3,3]$.

Range: $f(x)\geq 0$ and $f(0)=3$, so range $=[0,3]$.

\textbf{(b)} $g(x)=\frac{x-1}{x^2-4}=\frac{x-1}{(x-2)(x+2)}$: need $x\neq\pm 2$.

Domain $= \R\setminus\{-2,2\}$. Range: $g$ takes all real values except... at $x=1$, $g(1)=0$. Analysis shows range $=\R$ (every value achieved by continuity and the behaviour near the poles).

\textbf{(c)} $h(x)=\ln(x^2-1)$: need $x^2-1>0$, i.e.\ $|x|>1$.

Domain $=(-\infty,-1)\cup(1,\infty)$. As $|x|\to 1^+$, $\ln(x^2-1)\to-\infty$; as $|x|\to\infty$, $\ln(x^2-1)\to\infty$. Range $=\R=(-\infty,\infty)$.

\textbf{(d)} $k(x,y)=\sqrt{x}+\sqrt{y}$: need $x\geq 0$ and $y\geq 0$.

Domain $= \R^2_+ = [0,\infty)^2$. Range: $k(x,y)\geq 0$ always; for any $z\geq 0$, take $x=z^2$, $y=0$: $k(z^2,0)=z$. Range $=[0,\infty)$.
\end{solbox}

\begin{solbox}
\textbf{B2. Injectivity and Surjectivity}

\textbf{(a)} $f(x)=3x-7$: \textbf{Bijective.}

Injective: $3x_1-7=3x_2-7 \Rightarrow x_1=x_2$. \checkmark

Surjective: for any $y\in\R$, $x=(y+7)/3$ gives $f(x)=y$. \checkmark

\textbf{(b)} $g(x)=x^3-x=x(x-1)(x+1)$: \textbf{Neither.}

Not injective: $g(0)=0=g(1) \cdot$... let's check: $g(0)=0$, $g(-1)=(-1)(0) = 0$ wait: $g(-1)=(-1)(-2)(0)=0$. So $g(0)=g(-1)=0$ with $0\neq -1$. Not injective.

Not surjective? Actually it IS surjective since $\lim_{x\to\pm\infty}g(x)=\pm\infty$ and $g$ is continuous, so by IVT every value in $\R$ is achieved. So $g$ is \textbf{surjective but not injective}.

\textbf{(c)} $h:\R\to\R^+$, $h(x)=e^x$: \textbf{Bijective.}

Injective: $e^{x_1}=e^{x_2}\Rightarrow x_1=x_2$. \checkmark

Surjective (onto $\R^+$): for any $y>0$, $x=\ln y\in\R$ gives $h(x)=e^{\ln y}=y$. \checkmark

\textbf{(d)} $k:[0,\pi]\to[-1,1]$, $k(x)=\cos x$: \textbf{Bijective.}

Injective: $\cos$ is strictly decreasing on $[0,\pi]$, so injective. \checkmark

Surjective (onto $[-1,1]$): $k(0)=1$, $k(\pi)=-1$, continuous, so by IVT every value in $[-1,1]$ is achieved. \checkmark

\textbf{(e)} $\ell:\R_+\to\R_+$, $\ell(x)=x^3$: \textbf{Bijective.}

Injective: $x^3$ is strictly increasing on $\R_+$. \checkmark

Surjective: for any $y\geq 0$, $x=y^{1/3}\geq 0$ gives $\ell(x)=y$. \checkmark
\end{solbox}

\begin{solbox}
\textbf{B3. Inverse Functions}

\textbf{(a)} $f(x)=5x+3$: Set $y=5x+3 \Rightarrow x=\frac{y-3}{5}$.

So $f^{-1}(y)=\dfrac{y-3}{5}$, domain $\R$.

\textbf{(b)} $g(x)=\dfrac{x+1}{x-2}$, $x\neq 2$: Set $y=\dfrac{x+1}{x-2}$.

$y(x-2)=x+1 \Rightarrow xy-2y=x+1 \Rightarrow x(y-1)=2y+1 \Rightarrow x=\dfrac{2y+1}{y-1}$.

So $g^{-1}(y)=\dfrac{2y+1}{y-1}$, domain $\R\setminus\{1\}$.

\textbf{(c)} $h(x)=e^{2x-1}$: Set $y=e^{2x-1} \Rightarrow \ln y = 2x-1 \Rightarrow x=\dfrac{\ln y+1}{2}$.

So $h^{-1}(y)=\dfrac{\ln y+1}{2}$, domain $(0,\infty)$.

\textbf{(d)} $k(x)=x^3+2$ on $\R$: Set $y=x^3+2 \Rightarrow x^3=y-2 \Rightarrow x=(y-2)^{1/3}$.

So $k^{-1}(y)=(y-2)^{1/3}$, domain $\R$ (cube roots defined for all reals).
\end{solbox}

\begin{solbox}
\textbf{B4. Composition}

$f(x)=2x^2-1$, $g(x)=\sqrt{x+3}$ with domain $[-3,\infty)$.

\textbf{(a)} $(g\circ f)(x)=g(f(x))=\sqrt{(2x^2-1)+3}=\sqrt{2x^2+2}=\sqrt{2}\cdot\sqrt{x^2+1}$.

Domain: need $f(x)\in[-3,\infty)$, i.e.\ $2x^2-1\geq -3$, i.e.\ $x^2\geq -1$. This holds for all $x\in\R$.

Domain of $g\circ f$ is $\R$.

\textbf{(b)} $(f\circ g)(x)=f(g(x))=2(\sqrt{x+3})^2-1=2(x+3)-1=2x+5$.

Domain: need $x\in[-3,\infty)$ (the domain of $g$).

Domain of $f\circ g$ is $[-3,\infty)$.

\textbf{(c)} $(g\circ f)(2)=\sqrt{2(4)+2}=\sqrt{10}$

$(f\circ g)(1)=2(1)+5=7$

\textbf{(d)} $(g\circ f)(x)=\sqrt{2}\sqrt{x^2+1}$ vs.\ $(f\circ g)(x)=2x+5$. These are clearly different functions.

This illustrates that function composition is \textbf{not commutative}: $g\circ f\neq f\circ g$ in general.
\end{solbox}
\newpage
\begin{solbox}
\textbf{B5. Pre-image}

$f(x)=x^2-4$.

\textbf{(a)} $f^{-1}(\{0\})=\{x:x^2-4=0\}=\{x:x^2=4\}=\{-2,2\}$.

\textbf{(b)} $f^{-1}([-4,0])=\{x:-4\leq x^2-4\leq 0\}=\{x:0\leq x^2\leq 4\}=\{x:|x|\leq 2\}=[-2,2]$.

\textbf{(c)} $f^{-1}((0,5])=\{x:0<x^2-4\leq 5\}=\{x:4<x^2\leq 9\}=\{x:2<|x|\leq 3\}=[-3,-2)\cup(2,3]$.

\textbf{(d)} $f^{-1}(\{-5\})=\{x:x^2-4=-5\}=\{x:x^2=-1\}=\emptyset$.

(No real number has a negative square, so no input maps to $-5$.)
\end{solbox}

\begin{solbox}
\textbf{B6. Homogeneity}

\textbf{(a)} $f(K,L)=3K^{0.4}L^{0.6}$:

$f(tK,tL)=3(tK)^{0.4}(tL)^{0.6}=3t^{0.4}K^{0.4}t^{0.6}L^{0.6}=t^{1}\cdot 3K^{0.4}L^{0.6}=t^1 f(K,L)$.

\textbf{Homogeneous of degree 1} (CRS). By Euler's theorem: $K\cdot\frac{\partial f}{\partial K}+L\cdot\frac{\partial f}{\partial L}=1\cdot f(K,L)$.

\textbf{(b)} $g(x,y)=x^2+xy+y^2$:

$g(tx,ty)=(tx)^2+(tx)(ty)+(ty)^2=t^2x^2+t^2xy+t^2y^2=t^2(x^2+xy+y^2)=t^2g(x,y)$.

\textbf{Homogeneous of degree 2}.

\textbf{(c)} $h(x,y)=\dfrac{x^2y}{x^3+y^3}$:

$h(tx,ty)=\dfrac{(tx)^2(ty)}{(tx)^3+(ty)^3}=\dfrac{t^3 x^2 y}{t^3(x^3+y^3)}=\dfrac{x^2y}{x^3+y^3}=t^0\cdot h(x,y)$.

\textbf{Homogeneous of degree 0} (like Walrasian demand: scale-invariant).

\textbf{(d)} $k(x,y)=x+y+1$:

$k(tx,ty)=tx+ty+1=t(x+y)+1\neq t^r(x+y+1)$ for any $r$ (the constant 1 is unaffected by scaling).

\textbf{Not homogeneous} of any degree.
\end{solbox}

\begin{solbox}
\textbf{B7. Convexity and Concavity}

\textbf{(a)} $f(x)=x^4$: $f''(x)=12x^2\geq 0$ for all $x$, with equality only at $x=0$.

\textbf{Convex} (not strictly convex everywhere, since $f''(0)=0$, but it is strictly convex in the sense that $f(\lambda x+(1-\lambda)y)<\lambda f(x)+(1-\lambda)f(y)$ for $x\neq y$).

\textbf{(b)} $g(x)=x^3$: $g''(x)=6x$. $g''(x)>0$ for $x>0$ (convex) and $g''(x)<0$ for $x<0$ (concave). \textbf{Neither convex nor concave} on all of $\R$.

\textbf{(c)} $h(x)=-2x^2+4x-1$: $h''(x)=-4<0$ for all $x$. \textbf{Strictly concave.}

\textbf{(d)} $k(x)=|x|$: For $x\geq 0$, $k''=0$ (weakly convex); for $x\leq 0$, $k''=0$. Using the definition: $k(\lambda x+(1-\lambda)y)\leq\lambda k(x)+(1-\lambda)k(y)$ by the triangle inequality.

$|\lambda x+(1-\lambda)y|\leq |\lambda x|+|(1-\lambda)y|=\lambda|x|+(1-\lambda)|y|$ (since $\lambda,1-\lambda\geq 0$). \textbf{Convex} (weakly, not strictly).
\end{solbox}

\begin{solbox}
\textbf{B8. Monopolist Economic Application}

\textbf{(a)} Profit function:
\[
  \pi(Q) = \text{Revenue} - \text{Cost} = P(Q)\cdot Q - C(Q) = (100-2Q)Q - (Q^2+10Q+50)
\]
\[
  \pi(Q) = 100Q - 2Q^2 - Q^2 - 10Q - 50 = -3Q^2 + 90Q - 50
\]
Domain: $Q\geq 0$ and $P(Q)=100-2Q\geq 0 \Rightarrow Q\leq 50$. So domain $=[0,50]$.

\textbf{(b)} $\pi(Q)=-3Q^2+90Q-50$ is a downward parabola. Maximum at $Q^*=15$ (see (d)):
\[
  \pi_{\max} = -3(225)+90(15)-50 = -675+1350-50 = 625
\]
At endpoints: $\pi(0)=-50$, $\pi(50)=-3(2500)+4500-50=-50$.

Range $= [-50, 625]$.

\textbf{(c)} $\pi''(Q) = -6 < 0$ for all $Q$. Therefore $\pi$ is \textbf{strictly concave}.

Economic meaning: each additional unit of output raises profit by less and less (eventually reducing it). The profit-maximising $Q$ is uniquely determined and is a global maximum.

\textbf{(d)} First-order condition: $\pi'(Q)=-6Q+90=0 \Rightarrow Q^*=15$.

$\pi(15) = -3(225)+90(15)-50 = -675+1350-50 = \mathbf{625}$.

Maximum profit is \textbf{\$625} at output $Q^*=15$.

\textbf{(e)} Break-even: $\pi(Q)=0$:
\[
  -3Q^2+90Q-50=0 \Rightarrow 3Q^2-90Q+50=0
\]
\[
  Q = \frac{90\pm\sqrt{8100-600}}{6} = \frac{90\pm\sqrt{7500}}{6} = \frac{90\pm 50\sqrt{3}}{6} = 15\pm\frac{25\sqrt{3}}{3}
\]
\[
  Q_1 = 15 - \frac{25\sqrt{3}}{3}\approx 15-14.43=0.57, \qquad Q_2=15+\frac{25\sqrt{3}}{3}\approx 29.43
\]
So $\pi^{-1}(\{0\})\approx \{0.57,\; 29.43\}$.

\textbf{Interpretation}: The firm breaks even at approximately $Q=0.57$ (barely above shutdown) and $Q=29.43$ (the high-output break-even point). For $Q\in(0.57,\,29.43)$, the firm makes positive profit; outside this range, it makes a loss. The firm should produce $Q^*=15$ to maximise profit.
\end{solbox}






\end{document}